\subsection*{The Arithmetic Signature}
\label{subsec:arithmetic-signature}

\begin{definition}[Arithmetic Signature]
\label{def:arithmetic-signature}
The \textbf{arithmetic signature} is the collection of symbols used to view
number systems as structures: arithmetic operations such as $+$ and $\cdot$,
distinguished constants such as $0$ and $1$, and relations such as $\leq$ when
order is included.
\end{definition}

\begin{remark*}[Standard quantified statement]
\[
\mathcal{L}_{\mathrm{arith}}
=\{+,\cdot,0,1,\leq,\ldots\},
\]
with the visible symbols chosen according to the intended number-system
structure.
\end{remark*}

\begin{remark*}[Definition predicate reading]
The arithmetic signature is the symbol list used before a particular number
system interprets those symbols.
\end{remark*}

% Omitted Negated quantified statement: a signature is specified data, not a
% predicate with a useful mathematical negation in this section.

\begin{remark*}[Interpretation]
Different number systems interpret the same visible arithmetic symbols on
different carrier sets.  For example, addition on $\mathbb{N}$ and addition on
$\mathbb{R}$ have the same symbol but different domains and different
structural consequences.
\end{remark*}

\begin{dependencies}
\begin{itemize}
  \item \hyperref[def:first-order-signature]{First-Order Signature}
  \item \hyperref[def:unary-operation]{Unary Operation}
  \item \hyperref[def:binary-operation]{Binary Operation}
  \item \hyperref[def:nullary-operation]{Nullary Operation}
\end{itemize}
\end{dependencies}
